\section{Chapter~\ref{sec:chemoutput}. The chemical output}

The two wires to the heart and their different speeds, negative feedback in hormone cascades, the pulse-frequency code, and the daily oscillator and its locking to light are measured directly; the bound $K<8$ is a theorem of control theory derived in the chapter; pulsations generated by the cascade feedback itself are a hypothesis with strong experimental support.

{\footnotesize
\setlength{\tabcolsep}{3pt}\renewcommand{\arraystretch}{1.15}
\begin{longtable}{>{\raggedright}p{0.5cm}>{\raggedright}p{4.0cm}>{\raggedright}p{5.6cm}>{\raggedright}p{2.3cm}>{\raggedright\arraybackslash}p{1.7cm}}
\toprule
No. & Claim & Experiment and what was measured & Source & Status \\
\midrule\endfirsthead
\toprule
No. & Claim & Experiment and what was measured & Source & Status \\
\midrule\endhead
1 & The heart's pacemaker beats by itself about $100$ times a minute; at rest the brake is on, and the rate is usually about $60$--$70$ & Humans, complete block of both wires to the heart: the intrinsic rate is above the resting rate and declines with age, about $100$ beats per minute in adults. So the brake dominates at rest. The resting rate of $60$--$70$ is a typical value, not cited separately. & \cite{JoseCollison1970ihr} & measured \\
2 & The \nt{NORA} gas and the \nt{ACET} brake are low-pass filters, the brake is faster~\eqref{eq:gasbrake} & Dog, frequency-modulated stimulation of the vagal or sympathetic cardiac nerve and the transfer function to atrial rate: both pathways are low-pass filters; the sympathetic one has a much lower corner frequency plus a pure delay of about $1.7\,\text{s}$. The characteristics depend on the mean tone, so the linear formula~\eqref{eq:gasbrake} is an approximation. & \cite{Berger1989} & measured \\
3 & The brake is fast because \nt{ACET} opens the potassium channel without a second messenger, while \nt{NORA} acts through a chain of reactions & Membrane patches of atrial cells: the potassium channel opened by acetylcholine is opened by the $\beta\gamma$ subunits of the receptor-coupled G protein, with no second messenger inside the cell. The \nt{NORA} pathway via cAMP is not cited here. & \cite{Logothetis1987gbg} & measured \\
4 & Blood circulates through the body in about a minute; \nt{ADRE} from the bloodstream reaches every organ with a receptor & A general order-of-magnitude estimate; stated as such in the chapter, no source cited. & --- & estimate \\
5 & A hormone cascade is enclosed in negative feedback: the final hormone inhibits the first stages & Review of experiments: adrenocortical hormones suppress ACTH release by the pituitary and releasing-hormone release by the hypothalamus in three time domains: fast, intermediate, and slow. & \cite{KellerWood1984fb} & measured \\
6 & Three identical stages are stable for $K<8$~\eqref{eq:cascadestable}; at the boundary the period is $2\pi\tau/\sqrt3\approx3.6\tau$ & Derived in the chapter: at the frequency $\sqrt3/\tau$ three stages give a $180^\circ$ phase shift and an attenuation by $8$. & derived in chapter & theory \\
7 & The level error is $1/(1+K)$, so such a regulator cannot do better than about $10\,\%$ (Proposition~\ref{prop:cascade}) & A consequence of row~6 for proportional feedback. The real axis has feedback at several stages and in several time domains (row~5), so the bound applies to model~\eqref{eq:cascade} and is not measured. & derived in chapter & theory \\
8 & At $K=4$ and $7$ the level settles; at $K=12$, steady pulsations with period $3.7$--$3.9\,\tau$ (Fig.~\ref{fig:cascade}) & Computed with \texttt{models/\allowbreak chem\_\allowbreak models.py}: at $K=12$ successive periods are $3.9$, $3.7$, $3.8$, $3.7\,\tau$; at $K=4$ the swing at the end of the run is $0.003$. & \texttt{models/\allowbreak chem\_\allowbreak models.py} & model \\
9 & The stress hormone is released in pulses about once an hour; the pulses are generated by the feedback itself, without a separate oscillator & Rats, constant infusion of releasing hormone: ACTH and corticosterone oscillate at the physiological, nearly hourly frequency, so no rhythmic input from the hypothalamus is needed for the pulses. A pituitary--adrenal model with delayed feedback produces such oscillations. & \cite{Walker2012glc, Walker2010} & hypothesis \\
10 & The receiver reads the pulse frequency, not the level & Monkeys with no release of their own gonadotropin-releasing hormone: constant infusion does not restore the release of pituitary hormones, hourly pulses do; switching back to constant infusion silences the response again. Receptor fatigue as a high-pass filter is an interpretation. & \cite{Belchetz1978} & measured \\
11 & Different pulse frequencies act differently on different cells of one gland & The same monkeys: raising the rate to $2$--$5$ pulses per hour lowers both pituitary hormones; lowering it to one per $3$~hours lowers one (LH) and raises the other (FSH), so their ratio is set by the rate. & \cite{Wildt1981gnrh} & measured \\
12 & The daily rhythm is generated by intracellular negative feedback with a delay of several hours & Review: clock-gene proteins suppress the transcription of their own genes with a delay; the transcription--translation loop gives a period of about a day. & \cite{TakahashiJS2017} & measured \\
13 & The master oscillator is a group of tens of thousands of neurons that synchronizes the rest of the body & Review: the suprachiasmatic nucleus is the master daily oscillator; its individual neurons in culture are autonomous oscillators, and the network synchronizes them; in the mouse, about $10^4$ neurons on each side. & \cite{Welsh2010scn} & measured \\
14 & The intrinsic period in humans is $24.18$~hours & Humans in dim light on a schedule decoupled from the day: the period of the melatonin, temperature, and cortisol rhythms averages $24.18$~hours in young and older people, with a narrow spread. & \cite{Czeisler1999} & measured \\
15 & Light adjusts the oscillator like a phase-locked loop~\eqref{eq:pll}; a shift of about $11$~min per day is needed & Humans, a single $6.7$-hour bright-light exposure at different times of day: a type~1 phase response curve with a $5$-hour range, delay before the body-temperature minimum and advance after it. Equation~\eqref{eq:pll} is the standard model; $11$~min $=0.18$~hours is arithmetic. & \cite{Khalsa2003prc} & measured \\
16 & Without light there is no locking, and the rhythm drifts to its own period & Totally blind people with no light perception living in ordinary society: in about half of them the melatonin and cortisol rhythms run at their own period, which does not match the day. & \cite{Sack1992blind} & measured \\
\bottomrule
\end{longtable}}
